\section{Additional Diagnostic Results}
\label{app:additional_results}

This appendix provides supporting diagnostics for the main results in
Table~\ref{tab:main_results}. We keep the discussion brief: the purpose is to
make the reported evaluation auditable without shifting the paper's focus away
from the evaluate--adapt--simulate workflow.

\subsection{Evaluation Sizes}

Table~\ref{tab:dataset_sizes} reports the number of evaluated units per split.
Eval-B is counted in paired comparisons rather than individual scenarios,
because each prediction compares two matched records.

\begin{table}[h]
\centering
\small
\begin{tabular}{@{}lr@{}}
\toprule
\textbf{Split} & \textbf{Evaluation units} \\
\midrule
Eval-A rule identification & 1,512 \\
Eval-B comparative statics & 280 pairs \\
Placebo & 504 \\
Social OOD & 576 \\
Career OOD & 576 \\
\bottomrule
\end{tabular}
\caption{Evaluation sizes for the diagnostic stage.}
\label{tab:dataset_sizes}
\end{table}

\subsection{Parse Rates}

Table~\ref{tab:parse_rates} checks whether low accuracy could be explained by
formatting failures. With the exception of GPT-4 on Eval-B, parse rates are
near ceiling. The main conclusions are therefore about decision errors rather
than output-format noncompliance.

\begin{table*}[h]
\centering
\small
\setlength{\tabcolsep}{5pt}
\begin{tabular}{@{}lrrrrr@{}}
\toprule
\textbf{Model} & \textbf{Eval-A} & \textbf{Eval-B} & \textbf{Placebo} & \textbf{Social OOD} & \textbf{Career OOD} \\
\midrule
DeepSeek-R1-671B & 99.9 & 99.3 & 98.4 & 99.8 & 97.0 \\
Qwen & 99.5 & 100.0 & 99.4 & 99.8 & 99.7 \\
GPT-4 & 100.0 & 88.6 & 100.0 & 100.0 & 100.0 \\
Llama-3.1 & 100.0 & 100.0 & 100.0 & 100.0 & 100.0 \\
Llama-3 8B & 100.0 & 100.0 & 100.0 & 100.0 & 100.0 \\
Qwen2 7B & 100.0 & 100.0 & 100.0 & 100.0 & 100.0 \\
\bottomrule
\end{tabular}
\caption{Answer parse rates (\%). Unparsed responses are counted as incorrect in accuracy.}
\label{tab:parse_rates}
\end{table*}

\subsection{Eval-B by Perturbation Type}

Table~\ref{tab:evalb_perturbations} decomposes the comparative-static task by
which latent factor changes. Models are strongest when the peer action level or
weighted reference aggregate increases, but weaker when the perturbation depends
on comparison intensity or the redistribution of closeness weights. This
supports the interpretation that models have partial directional sensitivity but
do not fully internalize the structural rule.

\begin{table*}[h]
\centering
\small
\setlength{\tabcolsep}{4pt}
\begin{tabular}{@{}lrrrrr@{}}
\toprule
\textbf{Model} & \textbf{$F$ up} & \textbf{$\alpha_i$ up} & \textbf{Peer action up} & \textbf{$R_i$ up} & \textbf{Top weight up} \\
\midrule
DeepSeek-R1-671B & 51.8 & 51.8 & 73.2 & 75.0 & 46.4 \\
Qwen & 53.6 & 57.1 & 91.1 & 92.9 & 55.4 \\
GPT-4 & 51.8 & 44.6 & 75.0 & 89.3 & 51.8 \\
Llama-3.1 & 64.3 & 42.9 & 73.2 & 76.8 & 51.8 \\
Llama-3 8B & 53.6 & 55.4 & 67.9 & 66.1 & 48.2 \\
Qwen2 7B & 44.6 & 64.3 & 66.1 & 60.7 & 46.4 \\
\bottomrule
\end{tabular}
\caption{Eval-B accuracy (\%) by comparative-static perturbation.}
\label{tab:evalb_perturbations}
\end{table*}

\subsection{Social OOD by Match Distance}

Table~\ref{tab:social_distance} reports the social OOD split. The task does not
expose the numeric $b/c$ ratio and instead uses semantic benefit/cost tiers. The
mid bucket is consistently hardest, suggesting that models rely on coarse
ordering of semantic cues and struggle when the benefit-cost match is less
separable.

\begin{table*}[h]
\centering
\small
\setlength{\tabcolsep}{6pt}
\begin{tabular}{@{}lrrr@{}}
\toprule
\textbf{Model} & \textbf{Close} & \textbf{Mid} & \textbf{Far} \\
\midrule
DeepSeek-R1-671B & 58.3 & 50.9 & 63.9 \\
Qwen & 67.6 & 60.2 & 72.2 \\
GPT-4 & 74.1 & 60.6 & 75.7 \\
Llama-3.1 & 64.8 & 52.3 & 72.2 \\
Llama-3 8B & 59.7 & 52.8 & 70.1 \\
Qwen2 7B & 60.2 & 51.9 & 68.8 \\
\bottomrule
\end{tabular}
\caption{Social OOD accuracy (\%) by semantic benefit/cost match-distance bucket.}
\label{tab:social_distance}
\end{table*}

\subsection{Eval-A Rule Confusions}

Table~\ref{tab:evala_rule_confusions} expands the Eval-A error analysis. The
dominant mistakes are not random: models repeatedly replace the closeness-
weighted rule with single-peer anchoring, uniform aggregation, pure private
choice, or equal mixing. These categories define the main negative examples for
the targeted fine-tuning stage.

\begin{table*}[h]
\centering
\small
\setlength{\tabcolsep}{3pt}
\begin{tabular}{@{}lrrrrrr@{}}
\toprule
\textbf{Model} & \textbf{Top/closest} & \textbf{Uniform} & \textbf{Private} & \textbf{Median} & \textbf{Counter} & \textbf{Mix} \\
\midrule
DeepSeek-R1-671B & 26.7 & 21.2 & 8.2 & 18.4 & 8.3 & 17.1 \\
Qwen & 32.2 & 20.5 & 10.3 & 17.4 & 3.4 & 15.5 \\
GPT-4 & 48.5 & 3.6 & 18.4 & 6.7 & 0.5 & 22.3 \\
Llama-3.1 & 33.4 & 7.6 & 27.1 & 11.2 & 7.8 & 12.8 \\
Llama-3 8B & 32.5 & 6.4 & 27.9 & 13.1 & 6.1 & 14.0 \\
Qwen2 7B & 34.0 & 9.5 & 25.7 & 15.2 & 1.6 & 14.0 \\
\bottomrule
\end{tabular}
\caption{Eval-A wrong-answer distribution (\% of each model's Eval-A errors).}
\label{tab:evala_rule_confusions}
\end{table*}

\subsection{Career OOD by Status Sensitivity Bucket}

Table~\ref{tab:career_alpha2} stratifies career sorting by relative-status
sensitivity. Accuracy remains near chance across low, mid, and high buckets,
including for GPT-4. This suggests that the failure is not isolated to one
status-sensitivity regime; models generally struggle with the threshold tradeoff
between absolute salary and relative rank.

\begin{table*}[h]
\centering
\small
\setlength{\tabcolsep}{6pt}
\begin{tabular}{@{}lrrr@{}}
\toprule
\textbf{Model} & \textbf{Low $\alpha_{2i}$} & \textbf{Mid $\alpha_{2i}$} & \textbf{High $\alpha_{2i}$} \\
\midrule
DeepSeek-R1-671B & 43.8 & 45.8 & 47.9 \\
Qwen & 47.4 & 49.5 & 51.6 \\
GPT-4 & 50.0 & 50.0 & 50.0 \\
Llama-3.1 & 47.9 & 52.1 & 51.0 \\
Llama-3 8B & 50.5 & 51.6 & 53.1 \\
Qwen2 7B & 47.4 & 46.4 & 49.0 \\
\bottomrule
\end{tabular}
\caption{Career OOD accuracy (\%) by relative-status sensitivity bucket. Accuracy remains close to chance across buckets.}
\label{tab:career_alpha2}
\end{table*}
